This chapter adopts a data-oriented perspective on the standardization and interoperability challenges discussed in the previous chapter. It first argues for treating the data and software dimensions of a Digital Twin as separable, complementary design concerns, and examines how the two challenges just discussed relate to this separation. It then characterizes the data that Digital Twins produce and consume, discussing the requirements that follow from the diversity of the domains in which Digital Twins are deployed. Finally, building on this characterization, it examines the architectural and semantic foundations required to support such data across individual Digital Twins and Digital Twin Ecosystems.

\section{Divide et Impera}

Digital Twins, and the platforms that support them, can be described along two complementary dimensions: a \textit{software} dimension, concerned with the applications, services, and infrastructural components through which Digital Twins are deployed, executed, and made accessible; and a \textit{data} dimension, concerned with the information that Digital Twins collect, process, store, and exchange. These two dimensions are not independent, but modular and complementary. Every software component of a Digital Twin, from the simplest monitoring dashboard to the most sophisticated simulation service, ultimately operates on data: it acquires it, transforms it, and produces new data to be stored and eventually consumed elsewhere. In this sense, the software dimension of a Digital Twin can be regarded as a layer of applications built on top of the data dimension, rather than as an independent architectural concern.

This observation motivates a deliberate decomposition of the problem. Rather than addressing the software and data dimensions jointly, and thereby risking the reproduction of the same fragmentation already observed in existing Digital Twin implementations, this thesis separates the two concerns and concentrates on the data dimension. This choice is not intended to diminish the relevance of software architectures, which remain necessary to deploy, orchestrate, and expose Digital Twin functionalities. It reflects, instead, the observation that many of the interoperability difficulties discussed in the previous chapter originate not from the diversity of software implementations per se, but from the heterogeneity of the underlying data: its models, schemas, semantics, and governance \cite{acharya2024interoperability,david2024interoperability}. Two Digital Twins built on entirely different software stacks can still interoperate if they share a common understanding of the data they exchange, whereas two Digital Twins built on the same software stack may fail to interoperate if their underlying data representations are incompatible.

This separation can be made more concrete by revisiting the standardization and interoperability challenges discussed in \Cref{sec:challenges}. \Cref{tab:standardization_interoperability} decomposes each of these challenges, at both the individual Digital Twin and the Digital Twin Ecosystem level, into a software-related and a data-related component. In each cell, the software-related component concerns the mechanisms through which components are built or connected -- architectural building blocks, communication protocols, interfaces -- an area increasingly addressed by existing reference architectures and standardization efforts \cite{wang2022standards}. The data-related component concerns instead the models, representations, and shared meaning of the information involved, and is comparatively less mature, particularly at the Digital Twin Ecosystem level: standardizing a communication protocol does not by itself guarantee that the resulting data can be modelled consistently or interpreted uniformly across independently developed Digital Twins. This asymmetry motivates the choice, introduced above, to concentrate the remainder of this thesis on the data dimension.

\begin{table}[htbp]
\centering
\caption{Standardization and interoperability challenges at the Digital Twin (DT) and Digital Twin Ecosystem (DTE) levels, decomposed into their software and data components.}
\label{tab:standardization_interoperability}
\renewcommand{\arraystretch}{1.5} % Leggermente aumentato per far respirare il testo
\begin{tabular}{p{3.2cm} p{5.5cm} p{5.5cm}}
\toprule
\textbf{Challenge} & \textbf{DT} & \textbf{DTE} \\
\midrule
\textbf{Standardization} 
& 
\textit{Software:} architectural principles, building blocks. \newline
\textit{Data:} data models, transformation methodologies and metadata management.
& 
\textit{Software:} shared infrastructure, protocols, interfaces. \newline
\textit{Data:} common data models and schemas, DTP.
\\
\addlinespace
\textbf{Interoperability} 
& 
\textit{n/a}
& 
\textit{Software:} communication protocols, interface standards. \newline
\textit{Data:} common semantics, shared governance and metadata management.
\\
\bottomrule
\end{tabular}
\end{table}

Furthermore, the modularity between the data and software dimensions allows research to be conducted individually on each of them, without requiring the other to be addressed simultaneously. This thesis concentrates on the data dimension, while acknowledging that the software dimension remains a necessary complement to it.

Before this data dimension can be supported by an appropriate platform, however, it must first be understood on its own terms: what kinds of data Digital Twins actually produce and consume, and what requirements follow from their specific characteristics. The next section explores this question.

\section{Digital Twin Data}
\label{sec:dt_data}

Data occupies a distinctive position among the elements that make up a Digital Twin~\cite{tao2022}. Unlike the physical entity or the virtual model, which are specific to a given application, data is the element through which every other component of a Digital Twin operates and interacts: it is what the physical entity produces, what the virtual model consumes and updates, what services are built upon, and what is exchanged whenever two Digital Twins communicate. Liu et al. note that data plays a central role in constructing virtual models, establishing cyber-physical connections, and supporting the operation of a Digital Twin \cite{liu2022dtd}. At the same time, the properties that distinguish a Digital Twin from a conventional simulation or a Digital Shadow, discussed in \Cref{sec:history-and-definition}, impose requirements on this data that are not typically found in other data-intensive applications. Characterizing this data, before deciding how to manage it, amounts to a form of domain understanding: just as data-intensive projects typically begin by examining the domain and the data available before committing to a specific modelling or engineering approach, addressing Digital Twin data management must begin by examining what this data actually looks like.

\subsection{Requirements on Digital Twin Data}

From the perspective of its data and at a conceptual level, a Digital Twin can be described, in most cases, in quite simple terms: a physical entity/environment monitored through sensors, whose observations are processed by models or, increasingly, by machine learning/AI components, and whose outputs are in turn used to act on the physical entity. This description makes explicit why the software and data dimensions introduced above are complementary rather than independent: the models and services built on top of a Digital Twin are only as effective as the data made available to them, and their outputs are only as meaningful as they can be related back to the physical entity being monitored.

What varies substantially across applications is the specific type of data involved in this loop, and the requirements it imposes. In manufacturing, discussed in \Cref{sec:ref-arch} with reference to the five-dimensional Digital Twin model, this typically includes sensor measurements from machines and production lines alongside geometric models of the manufactured parts \cite{tao2022}. In the automotive domain, Deng et al. review Digital Twin applications spanning vehicle design, production, transportation, and battery management, in which data ranges from structured diagnostic and telemetry signals, such as those used for battery state monitoring, to the higher-frequency data required for real-time production monitoring and control \cite{deng2023automotive}. In precision agriculture, Zhang et al. describe Digital Twins that combine environmental and soil sensor data with remote sensing imagery to support decisions over considerably longer, biologically determined time scales, and report that fragmented data sources and inconsistent data formats remain a significant obstacle to their integration \cite{zhang2025agriculture}.

These examples illustrate that the requirements imposed by Digital Twin data are extremely shaped by the domain of application. Liu et al. observe more generally that Digital Twin data also exhibits the volume, variety, and velocity commonly associated with big data \cite{liu2022dtd}: volume, because Digital Twins are typically built around densely instrumented entities that generate data continuously; variety, because of the heterogeneity of sources, domains, and data types illustrated above; and velocity, because different parts of the same Digital Twin may require different processing latencies, from near-real-time control loops to periodically updated maintenance records. A further concern, raised in the agricultural example above and identified by Ouedraogo et al. as one of the central challenges in Digital Twin data management, is data quality: sensor-originated data is subject to noise, faults, and calibration issues, which existing data management architectures and techniques need to account for explicitly \cite{ouedraogo2025dtdm}.

These characteristics motivate the more specific requirements identified in the literature. Liu et al. translate them into concrete requirements on Digital Twin data gathering, mining, fusion, interaction, iterative optimization, and on-demand usage \cite{liu2022dtd}.

\subsection{Categorizing Digital Twin Data}

Beyond its heterogeneity, Digital Twin data can be organized along two complementary dimensions found in the literature, each of which has direct consequences on how such data should be modelled and managed.

A first, widely adopted distinction, introduced by the ISO 23247-3 standard for the digital representation of manufacturing elements, separates \emph{static information}, which does not change over the course of the entity's operation, such as its identification or nominal characteristics, from \emph{dynamic information}, which changes as the physical entity operates, such as its status, location, or the current shape of a piece of material undergoing a manufacturing process \cite{liu2022dtd}. Establishing this distinction for a given Digital Twin is itself an exercise in domain understanding, since it clarifies, ahead of any implementation choice, which parts of the data behave as reference or master data, typically low in volume and velocity and well suited to conventional, schema-based representations, and which parts behave as high-velocity, continuously generated observations, more naturally modelled as time-series or event data and requiring correspondingly different storage and processing mechanisms. In this sense, this dynamic distinction anticipates some of the data modelling requirements addressed in the remainder of this thesis.

A second, complementary categorization organizes Digital Twin data according to where, within a Digital Twin, it originates. Liu et al. distinguish six such categories: \emph{physical entity-related data}, such as sensor measurements describing equipment, materials, or workers; \emph{virtual model-related data}, such as the parameters, processes, and outcomes of simulation models; \emph{service-related data}, produced by the algorithms and analytics built on top of the Digital Twin; \emph{domain knowledge}, encoding rules, standards, and expert experience; \emph{connection data}, which relates corresponding pieces of physical entity-related and virtual model-related data, such as a real, sensor-measured temperature and the corresponding value simulated by a virtual model; and \emph{fusion data}, obtained by combining the previous categories into a more complete representation of the physical entity \cite{liu2022dtd}.

From a data engineering perspective, these six categories can be read as occupying different layers of a data architecture, echoing the classic distinction between data, information, and knowledge \cite{ackoff1989}. Physical entity-related and virtual model-related data sit at the \emph{raw data} layer: they are the raw observations and simulation outputs produced, respectively, by the physical and the virtual side of the Digital Twin. Connection and fusion data sit one layer above, at the \emph{reconciled} layer: they combine and integrate these raw observations into a coherent, contextualized representation of the physical entity, with fusion data in particular constituting the integration layer at which a Digital Twin's overall view of its physical counterpart is assembled. Service-related data and domain knowledge sit at the \emph{knowledge} layer: the former encodes the outputs of analytics and algorithms applied to the underlying information, while the latter encodes rules and expertise external to any single observation; together, they are what ultimately supports the decisions and actions a Digital Twin takes on its physical entity.

This layered reading has a direct consequence for the data models discussed later in this chapter. A data model intended to support meaningful interoperability across Digital Twins cannot be designed solely around raw, entity- or model-related observations, since these are, by construction, specific to each Digital Twin's own sources. It must instead operate at the information layer, where connection and fusion data already provide a reconciled representation of the physical entity, largely independent of how that entity happens to be sensed or simulated.

Ouedraogo et al. observe that the specific data management considerations arising from this heterogeneity differ across application domains, such as manufacturing, smart cities, and healthcare \cite{ouedraogo2025dtdm}, reinforcing the earlier observation that the variety of Digital Twin data is compounded, rather than absorbed, by the variety of domains in which Digital Twins are deployed.

Taken together, these categorizations show that even a single Digital Twin already involves several qualitatively different types of data, produced and consumed by different components, at different rates, and for different purposes, and that its origin, its layer within the data-information-knowledge architecture just outlined, and its position along the static/dynamic spectrum all carry direct implications for how it should be modelled. This internal heterogeneity is compounded, at the level of a Digital Twin Ecosystem, by the differences in data models, schemas, and semantics adopted by independently developed Digital Twins, discussed in \Cref{sec:challenges}. Characterizing Digital Twin data along the dimensions introduced in this section therefore provides the domain understanding necessary not only for designing the data management capabilities of an individual Digital Twin, but also for identifying the requirements that must be satisfied to support Digital Twin Ecosystems more broadly.

\section{Digital Twin Platform}

A Digital Twin Platform (DTP) can be defined as a software environment that provides the data management mechanisms in terms of infrastructure, data models and semantics required to host and support Digital Twins throughout their lifecycle, that hosts and supports Digital Twins throughout their lifecycle, from deployment to decommissioning. What distinguishes the perspective adopted in this thesis is not whether a DTP is a software environment, but what it is organized around: rather than being structured around the specific software modules and services that make up the Digital Twins it hosts, it is structured around the data those Digital Twins produce and consume. 

This distinction has a direct practical consequence. A Digital Twin Platform designed in this way is largely independent of the internal architecture of the Digital Twins deployed on top of it: the platform does not need to know how a given Digital Twin is implemented internally, only that the data it produces and consumes conforms to the platform's data models and is semantic representation standards. A Digital Twin built according to these data models and semantic standards will run on the platform and be able to interoperate with other Digital Twins built the same way, regardless of the specific software stack each of them uses internally. Standardizing on data models and data representations, rather than on software architectures, is therefore the foundation on which this thesis proposes to standardize Digital Twins; the software components of a Digital Twin, whatever their specific architecture, are consequently treated as applications that operate on top of a shared data foundation, rather than as the primary object of standardization.

A Digital Twin Platform, in this sense, does not require every Digital Twin to follow an identical architecture. Its role is instead to provide common capabilities and abstractions that reduce fragmentation and facilitate the construction and management of heterogeneous Digital Twin environments, representing a technological foundation upon which different Digital Twins can be built and, ultimately, connected to one another.

This characterization does not contradict the reference architectures discussed previously, such as the five-dimensional Digital Twin model, in which data already occupies a central, connective role among the physical entity, virtual model, connection, and service dimensions \cite{qi2021,tao2022}. Rather, it makes this centrality explicit at the level of the platform, and uses it as an organizing principle for the remainder of the thesis. Similarly, this perspective is consistent with service-oriented reference models such as Digital Twin as a Service, in which Digital Twin functionality is exposed through composable, standardized building blocks \cite{AHELEROFF2021101225}; in the perspective adopted here, however, the building blocks that are standardized in the first instance are those governing the underlying data, with services constructed on top of them.

Treating data as a first-class citizen in this sense means that the requirements and characteristics of Digital Twin data, discussed in \Cref{sec:dt_data}, are taken as the primary drivers of platform design, rather than as a consequence of decisions made at the software or application level. In many existing systems, data models, storage solutions, and integration mechanisms are chosen to satisfy the immediate needs of a specific application, and are only subsequently, if at all, reconciled with the requirements of other applications or Digital Twins operating on related data. This application-centric approach tends to reproduce, at the level of the platform, the same fragmentation already observed among independently developed Digital Twins. The perspective adopted in this thesis inverts this relationship: the characteristics of Digital Twin data are analysed independently of any specific application, and are used to define the data management capabilities that a Digital Twin Platform should provide, with software applications, including the Digital Twins themselves, developed on top of these capabilities rather than dictating them.

The remainder of this chapter details two complementary aspects of this data-oriented view of the Digital Twin Platform: the architectural and infrastructural foundations required to support Digital Twin data, referred to as data platforms; and the data models and governance mechanisms required to give this data a shared representation and meaning across heterogeneous Digital Twins.

\subsection{Data Platforms}

The first architectural consequence of treating data as a first-class citizen concerns the infrastructural foundations needed to acquire, store, process, and serve Digital Twin data at scale. This thesis refers to these foundations as \textit{data platforms}. A data platform is a centralized infrastructure that facilitates the ingestion, storage, management, and exploitation of large volumes of heterogeneous data.
It provides a collection of independent and composable services meeting the end-to-end needs of data pipelines, where:
\textit{centralized} means that a data platform is conceptually a single and unified component;
\textit{independent} means that changes in the implementation of a service do not affect other services;
\textit{composable} means that services have interfaces that enable an easy and frictionless composition;
and \textit{end-to-end} means that services cover the entire data life cycle.
Data platforms foster collaboration and shared governance (being centralized, data is unified following some integration and it is easier to ensure compliance with data protection and privacy laws through shared security and access control), and scalability (being implemented on a distributed infrastructure, it is easy to add storage and computing resources as needed)~\cite{francia2025process}. 


This definition is well suited to the Digital Twin setting. As characterized in \Cref{sec:dt_data}, DT data combines high volume with uneven velocity, non-trivial data quality issues, and a wide variety of sources and structures shaped by the application domain, and imposes requirements, such as universality and fusion, that go beyond those of more conventional, application-specific data. Adopting a data platform as the infrastructural foundation of a Digital Twin, or of a Digital Twin Ecosystem, means standardizing how such data is ingested, stored, and made available to the applications and services that depend on it, independently of their specific implementation, rather than requiring every Digital Twin to design and operate its own bespoke data infrastructure.

It is important to be precise about the scope of this contribution. The infrastructural component of a Digital Twin Platform can be understood as a data platform tailored to the characteristics of Digital Twin data described in \Cref{sec:dt_data}; realizing such a tailored platform is, however, a design and engineering effort in its own right, which is addressed directly in the Architectures part of this thesis, rather than something a generic data platform provides out of the box. Once instantiated, such a platform would primarily address the standardization challenge identified in the previous chapter: at the level of an individual Digital Twin, by offering common building blocks for the acquisition, storage, and management of its data; and at the level of a Digital Twin Ecosystem, by offering a shared infrastructure on which multiple, independently developed Digital Twins can be deployed. Even then, a data platform does not by itself address the problem of finding a meaningful and shared representation of the integrated data. Standardizing the technical substrate through which data flows does not guarantee that the data flowing through it is structured, described, or interpreted consistently across the Digital Twins that produce and consume it. Two Digital Twins deployed on the same data platform can still fail to interoperate, at the syntactic or semantic levels discussed in \Cref{sec:challenges}, if they rely on incompatible data models or assign different meanings to the same underlying concepts \cite{acharya2024interoperability,david2024interoperability}. Addressing this residual problem requires moving from the infrastructural to the semantic level, through the definition of data models and governance mechanisms, discussed next.

The requirements that a data platform must satisfy in the context of Digital Twin applications, and the extent to which existing data management technologies can address them, are examined in detail in the remainder of the thesis, together with a concrete instantiation developed in the context of precision agriculture.

\subsection{Data Models and Governance}

As anticipated above, standardizing the infrastructural layer through which data flows is a necessary, but not sufficient, condition for building an integrated, consistent and shared representation of the heterogeneous collected data. As discussed in the previous chapter, two systems can be technically, or even syntactically, connected without being able to correctly interpret the information they exchange \cite{kritzinger2018,fuller2020,acharya2024interoperability}. Addressing this problem requires the definition of data models and governance mechanisms that operate on top of the underlying data platform.

Before considering their specific role within a Digital Twin Platform, it is useful to recall briefly what these two notions mean in more general terms. A \textit{data model} specifies how a portion of reality is represented in digital form: which entities and properties are captured, how they relate to one another, and how they are structured and encoded. Data models vary considerably in structure, ridigity and expressiveness, from a relational schema listing tables, columns, and key relationships, to document-oriented models such as JSON schema specifying the structure of an exchanged document, to more connection-oriented models such as the property and knowledge graphs data models, the latter introducing the concept of a formal ontology that additionally captures the meaning of the represented concepts and the relationships holding among them, independently of any specific storage technology. \textit{Data governance}~\cite{Khatri2010}, in turn, denotes the set of policies, processes, and responsibilities put in place to ensure that data is managed consistently for its intended uses throughout its lifecycle; in practice, this typically includes maintaining metadata describing the data's meaning, provenance, tracking how data is transformed as it moves through a system, assigning ownership and access permissions, and managing how data models are allowed to evolve without breaking the applications that depend on them.

At the level of an individual Digital Twin, data models formalize the structure and meaning of the information associated with the represented physical entity, building on the categorization introduced in \Cref{sec:dt_data} to provide a shared reference for the components and services that operate on that data. At the level of a Digital Twin Ecosystem, this requirement extends to the definition of data models capable of accommodating data originating from multiple, independently developed Digital Twins, so that information produced by one Digital Twin can be correctly interpreted and reused by another. Governance mechanisms complement these data models by managing their evolution, documenting their metadata, and ensuring that the meaning associated with shared concepts, such as a measured physical quantity, remains consistent across the Digital Twins that produce and consume it.

Some steps in this direction have already been taken in the broader data platform literature, even outside the specific context of Digital Twins. Francia et al., for instance, propose MOSES, a technology-agnostic framework for metadata handling in big data platforms, built around a metadata repository that stores information about the data objects circulating in the platform and the processes that transform them \cite{francia2021moses}. Frameworks of this kind illustrate how governance mechanisms operating on metadata can be layered on top of a data platform to keep track of the meaning, provenance, and transformations of the underlying data.

Taken together, data models and governance mechanisms provide the semantic counterpart to the infrastructural standardization discussed above, and constitute the basis on which independently developed Digital Twins can be said to interpret exchanged information consistently.

\begin{partsummary}

The discussion in this chapter provides the foundation for positioning the research presented in the remainder of this thesis. Achieving interoperability across independently developed Digital Twins is commonly regarded, in both the research literature and in large-scale initiatives such as the UK's National Digital Twin programme, as one of the more advanced stages in the maturity of Digital Twin deployments \cite{klar2024maturity,dtarch}. As this chapter has already argued, this maturity is primarily a property of the data a Digital Twin exchanges, rather than of the software that produces or consumes it: two Digital Twins built on incompatible software stacks can still interoperate if they share a common understanding of their data, while two Digital Twins sharing the same software stack may fail to interoperate if their data representations diverge. This thesis adopts the perspective that interoperability, understood in this sense, is a desirable objective, but one that cannot be pursued directly without first standardizing the data dimension on which it depends.

Before independently developed Digital Twins can interoperate, it is necessary to establish what constitutes a Digital Twin and, in particular, how its data is structured, represented, managed, and exchanged. Doing so requires addressing the software and data dimensions separately. While both are necessary to fully characterize a Digital Twin, this thesis deliberately focuses on the data dimension, precisely the dimension on which interoperability chiefly depends, leaving the standardization of software architectures, interfaces, and services outside its scope.

Within this scope, standardization proceeds through a sequence of steps that mirrors well-established data-intensive methodologies such as CRISP-DM. Before any infrastructure or representation can be designed, it is first necessary to understand the data itself, its sources, its heterogeneity, and the requirements it imposes, as briefly introduced in \Cref{sec:dt_data}. Only once this understanding is in place does it becomes possible to design the infrastructure needed to collect and manage that data, and subsequently to define a coherent and efficient representation for it.

This thesis carries out this progression at the level of an individual Digital Twin. Once Digital Twin data is characterized, it addresses the infrastructural step through the investigation of data platforms as the foundation for acquiring, storing, processing, and making Digital Twin data available, and the representational step through the definition of data models and associated governance mechanisms capable of providing a consistent representation of the information produced and consumed within the Digital Twin.

\end{partsummary}